%%%%%%%%%%%%%%%%%%%%%%%%%%%%%%%%%%%%%%%%%%%%%%%%%%%%%%%%%%%%%%%%%%%%%%%%%%%%%%%
% Chapter 3: Tecnologías
%%%%%%%%%%%%%%%%%%%%%%%%%%%%%%%%%%%%%%%%%%%%%%%%%%%%%%%%%%%%%%%%%%%%%%%%%%%%%%%

%++++++++++++++++++++++++++++++++++++++++++++++++++++++++++++++++++++++++++++++
\section{Lenguaje para la lógica de la aplicación}
\label{3:sec1}

Por normal general cuando se habla de programación web en el lado del cliente, se tiende a pensar
de forma inmediata en Javascript, y en general, este razonamiento es totalmente válido. Pero dado que SIMDE
es una aplicación fuertemente orientada a objetos y con una gran base de código, se han valorado múltiples opciones.

\subsection{Coffescript}

\subsection{Dart}

Dart es un lenguaje de código abierto desarrollado por Google que permite desarrollador aplicaciones web, móvil, 
de servidor y también se puede utilizar en el \textit{Internet of Things}. 

\bigskip
Se ha considerado en el desarrollo de esta aplicación porque es un lenguaje orientado a objetos que utiliza una 
sintaxis similar a C\#. Además, aunque Google Chrome tiene una máquina virtual nativa para este lenguaje, es posible
transpilar el código a Javascript para los navegadores que no tiene este soporte nativo.

\begin{figure}[!th]
\begin{center}
\includegraphics[width=0.5\textwidth]{images/cap3/polymerlogo.eps}
\caption{Ejemplo}
\label{fig:PolymerLogo}
\end{center}
\end{figure}

Tras razonar detenidamente, Dart ha quedado descartado por diversas razones:

\begin{enumerate}
\item no va
\end{enumerate} 

\subsection{Typescript}

Typescript es un lenguaje libre y de código abierto desarrollado por Microsoft que actúa como un superconjunto
de Javascript, es decir incorpora los distintos estándares: ECMA5, ECMA6, ECMA7... Y además, como característica
destacable, añade comprobación de tipos en tiempo de compilación.

\bigskip
Este tipado no se refleja en el código final, de hecho una interfaz, por ejemplo,
añade 0 sobrecarga en el código final. Pero si que es interesante por las capacidades de
autocompletar (a través de Microsoft Intellisense)  que añade. Typescript se puede transpilar 
directamente a código javascript es5, el cuál es el estándar actual en todos los navegadores.    

\bigskip
Si tuviéramos que barajar una posible alternativa, sin duda la más destacada sería Flow, 
un comprobador de tipos para javascript desarrollado por Facebook. 

\bigskip
Pero existen múltiples razones para que haya escogido typescript sobre Flow:

\begin{enumerate}

\item Flow no está siendo parte de un proyecto de gran envergadura, se suponía que
 se introduciría en el desarrollo de react. Sin embargo no ha sido así.

\item Typescript es la opción de defacto para Angular 2, ha ido ganando mucho entusiasmo
 por parte de la comunidad y tiene un gran soporte.

\item Tengo cierta experiencia con Typescript.
\end{enumerate}

%++++++++++++++++++++++++++++++++++++++++++++++++++++++++++++++++++++++++++++++
\section{Tecnología para la integración modelo vista}
\label{3:sec2}

\subsection{Webcomponents}

Desde el inicio del proyecto, tenía claro que alguna librería se encargaría de realizar 
por mí el tedioso proceso de manipular el DOM. Actualmente existen múltiples librerías 
y frameworks que podían servirme en esta tarea, pero muchos de ellos (como por ejemplo Angular),
 no son lo suficientemente flexibles y acaban condicionando el desarrollo de la aplicación.

\bigskip
Lo que SIMDE necesitaba era un proyecto que hiciera uso de los Web components, 
una serie de apis que permiten crear etiquetas html personalizadas y reutilizables. 
Los webs componentes consisten en 4 características principales que pueden usarse 
por separado o todos juntos:

\begin{enumerate}

\item \textbf{Elementos personalizados.}
\item \textbf{Shadow DOM.}
\item \textbf{HTML Imports.}
\item \textbf{HTML Templates.}

\end{enumerate}

\subsection{Polymer}
Probablemente la librería basada en Web Componentes más conocida es Polymer,
una librería desarrollada por Google y anunciada en 2013. 

\bigskip

Pero a pesar de que las múltiples ventajas de Polymer existe una librería cuya comunidad es espléndida,
 y no es otra que React.

\subsection{React}
Existen una gran cantidad de motivos para escoger React sobre Polymer: Como por ejemplo, 
que ahora mismo se utiliza en decenas de aplicaciones importantes, como netflix, airbnb, Wallmart… (TODO INCLUIR CITA)

\bigskip
Que la comunidad es impresionantemente activa y cuenta con una gran cantidad de usuarios dispuestos a ayudar, así como cuenta con muchísima documentación y muchísimas implementaciones de librerías de terceros.

\bigskip
React utiliza un híbrido entre html y javascript denominado jsx, como también tiene soporte para 
Typescript, en este caso utilizamos tsx, y se basa en un “unidirectonial data-flow”. 

\bigskip
Ademas React implementa operaciones sobre el DOM virtual de tal forma que las operaciones sobre
el verdadero DOM sean eficientes.

%++++++++++++++++++++++++++++++++++++++++++++++++++++++++++++++++++++++++++++++
\section{Tecnología para hacer la build}
\label{3:sec3}

\subsection{Gulp/Grunt}

\subsection{Webpack}
\bigskip
Una de las grandes herramientas de 2016 que acabó por cambiar el flujo de muchos 
desarrolladores web y desbancó a tasks runners como Grunt y Gulp fue webpack.

\bigskip
Para poder integrar todo este código y resolver el problema de los múltiples imports 
era necesario utilizar algún tipo de herramienta de gestión de paquetes, como por ejemplo 
commonjs, o requirejs. Sin embargo, webpack se encarga de resolver todas estas 
dependencias y crear statics assets para el navegador.

\bigskip
Uno de los mayores puntos a favor es que es altamente configurable, existen muchísimos 
plugins de webpack que permiten hacer preprocesamiento de css, tratamiento de imágenes,
minimización codigo, étctera.. 

\bigskip
Actualmente en la aplicación de SIMDE webpack  se encarga de compilar el código de 
typescript, el código de react y de generar un bundle 

%++++++++++++++++++++++++++++++++++++++++++++++++++++++++++++++++++++++++++++++
\section{Tecnología para la documentación}
\label{3:sec4}

Para integrar la documentación en la nueva aplicación web de SIMDE resultaba obvio que esta documentación
estuviera también en formato web. Para esto existían muchas alternativas, desde un conjunto de ficheros
html hasta un pequeño cms. 

Dado que la documentación es bastante extensa pero que en realidad, no es mas que un documento 
que se redactará en una ocasión y se le irán realizando pequeñas ampliaciones y/o correcciones
se optó por una solución diferente, los generadores de contenido estático.

\subsection{Generadores de contenido estático}

Los gneeradores de contenido estático 

\subsection{Hexo}



